\section{Concept Graph Attention}
\label{sec:main-cga}

CGA is an attention layer whose memory is a fixed bank of concept identities
rather than a table of token keys and values.
The ``graph'' is the runtime bipartite interaction between tokens and concept
nodes, with soft read and write weights computed by attention.
Each layer has $C$ concept slots and $H$ heads.
For head $h$, slot $j$ has a fixed unit address
\[
a_{hj}\in \mathbb{R}^{d_a},
\qquad
\|a_{hj}\|_2=1 .
\]
The addresses are initialized once and remain fixed; they are concept
identities, not online sequence state.
The online state for a sequence is the value memory
\[
M_b\in\mathbb{R}^{H\times C\times d_v^c},
\]
where $b$ indexes a causal processing block and $d_v^c$ is the per-head concept
value width.
Thus the persistent runtime state per layer is
\[
S_{\rm CGA}=C H d_v^c .
\]
This is the main state-size difference from token attention: attention stores a
new key-value entry for every token, while CGA stores a bounded value table over
fixed concept nodes.

\paragraph{Projection and normalization.}
Given hidden states $x_t\in\mathbb{R}^{d_{\rm model}}$, CGA forms normalized
read queries and write keys, concept values, write gates, and retention rates:
\[
q_{th}=\operatorname{normalize}(Q_h x_t),\qquad
k_{th}=\operatorname{normalize}(K_h x_t),
\]
\[
v_{th}=V_h x_t,\qquad
g_t=\sigma(w_g^\top x_t+b_g).
\]
The query and key maps learn how tokens address the fixed concept identities.
The value map learns what payload should be stored in concept memory, and the
write gate learns how much a token should modify memory.

\paragraph{Dense concept read.}
CGA reads by applying the standard attention softmax over concept slots rather
than over prefix tokens.
Let $\gamma>0$ be a learned read scale.
For token $t$ in block $b$,
\[
\alpha_{thj}
=
\frac{\exp(\gamma q_{th}^{\top}a_{hj})}
{\sum_{r=1}^{C}\exp(\gamma q_{th}^{\top}a_{hr})},
\qquad
z_{th}=\sum_{j=1}^{C}\alpha_{thj}M_{b,hj}.
\]
The multi-head output is
\[
y_t=O[z_{t1};\ldots;z_{tH}].
\]
This is ordinary differentiable attention except that the keys are fixed
concept identities and the values are runtime concept memories.

\paragraph{Continuous concept write.}
CGA uses a read-before-write protocol.
The tokens in a block first read the memory state $M_b$; after those outputs
are emitted, the block writes its values into memory to produce $M_{b+1}$.
For a write delay $\delta\ge0$, define
\[
\bar k_{th}=k_{t-\delta,h},
\]
with invalid delayed positions masked out.
MQAR uses $\delta=1$ so that a key token can bind to the following value token;
general language modeling can use $\delta=0$.
The write distribution is another concept softmax:
\[
\beta_{thj}
=
\frac{\exp(\gamma \bar k_{th}^{\top}a_{hj})}
{\sum_{r=1}^{C}\exp(\gamma \bar k_{th}^{\top}a_{hr})}.
\]
For block $b=[s,e)$, CGA updates each concept value by
\[
M_{b+1,hj}
=
\rho_{bh} M_{b,hj}
+
\sum_{t=s}^{e-1}
\mathbf{1}_{t\ge \delta}\,
g_t\,\beta_{thj}\,v_{th}.
\]
The write is continuous: every valid token can write to every concept with a
different weight, and the soft weights make the write path differentiable
during training.

\paragraph{Trainable retention.}
The factor $\rho_{bh}\in(0,1]$ controls how much old concept memory survives
when a new block is written.
CGA computes tokenwise positive decay rates $\lambda_{th}$ and accumulates them
over the block:
\[
\rho_{bh}
=
\exp\left(-\sum_{t=s}^{e-1}\lambda_{th}\right).
\]
Our default parameterization uses a PoST-style ordered spectrum for the base
decay rates:
\[
A_h=A_0+\sum_{i<h}\operatorname{softplus}(\Delta_i),
\]
combined with a learned input-dependent step size and optional
position-adaptive scaling.
This gives CGA the same kind of trainable forgetting mechanism used by modern
recurrent mixers, but the recurrence acts on concept value memory rather than
on token-indexed state.

\paragraph{Training and inference.}
During training, CGA is optimized end-to-end with the task loss.
Gradients flow through the read weights $\alpha$, the write weights $\beta$,
the written values $v_t$, the write gate $g_t$, and the retention factors
$\rho_{bh}$.
The concept addresses $a_{hj}$ are fixed identities, so the model learns to use
them by changing the query, key, value, gate, retention, and output
projections.
During inference, the learned parameters are fixed.
The concept memory $M_b$ is still updated by the same forward rule, so future
tokens condition on information written from the observed prefix.
This is the sense in which CGA performs test-time learning: it learns the
contents of a bounded memory during the forward pass, without running
backpropagation at inference time.

\paragraph{Relation to ordinary attention.}
Standard causal attention computes
\[
\operatorname{Attn}(X)_t
=
\sum_{s<t}
\operatorname{softmax}_s(q_t^\top K_s)V_s,
\]
where the memory basis is the set of prefix token positions.
CGA computes
\[
\operatorname{CGA}(X)_t
=
\sum_{j=1}^{C}
\operatorname{softmax}_j(q_t^\top a_j)M_{tj},
\]
where the memory basis is a fixed set of concept identities and only the value
contents $M_{tj}$ evolve over the prefix.
Thus CGA is not a different loss or optimizer around attention.
It is an attention replacement that moves the read/write memory from tokens to
concepts.

\begin{algorithm}[t]
\caption{\textsc{ConceptGraphAttentionForward}}
\label{alg:main-cga}
\begin{algorithmic}[1]
\Statex Input: hidden states $x_{1:T}$, fixed concept addresses
$a_{1:C}$, block size $B$.
\State Initialize value memory $M_0\gets0$.
\For{causal blocks $b=[s,e)$}
    \State Project block queries $q$, write keys $k$, values $v$, gates $g$,
    and retention rates $\lambda$.
    \State Read $M_b$ with dense concept attention over fixed addresses.
    \State Emit block outputs through the output projection.
    \State Form delayed write keys $\bar k$ and write weights $\beta$.
    \State Compute block retention
    $\rho_b=\exp(-\sum_{t=s}^{e-1}\lambda_t)$.
    \State Update
    $M_{b+1}\gets \rho_b M_b+\sum_{t=s}^{e-1}g_t\,\beta_t\,v_t$.
\EndFor
\State Return projected concept readouts.
\end{algorithmic}
\end{algorithm}
